        Let $y_i^{[j]}$ denote the response observed at time $t_i$ in trial $j$,
    with $i=0,\ldots,T-1$ and $j=1,\ldots,J$. In this experiment $T=15$ and
    $J=3$. The time points run from $0$ to $168$ in the original time units.
    The observed index set is
    $$
    \Omega = \{(i,j): y_i^{[j]}\ \hbox{is observed}\},
    \qquad n = |\Omega| = 45.
    $$

        The Neural ODE and Neural SDE were trained on normalized time and
    normalized response:
    $$
        \tau_i = \frac{t_i-t_{\min}}{t_{\max}-t_{\min}},
        \qquad
        z_i^{[j]} = \frac{y_i^{[j]}-\bar y}{s_y},
    $$
    where
    $$
        \bar y = \frac{1}{n}\sum_{(i,j)\in\Omega} y_i^{[j]},
        \qquad
        s_y =
            \left[
                \frac{1}{n}\sum_{(i,j)\in\Omega}(y_i^{[j]}-\bar y)^2
            \right]^{1/2}.
    $$
    Thus $\tau_i\in[0,1]$. The logistic PINN models also use normalized time,
    but their neural trajectories are expressed directly on the original
    response scale. Their data and physics losses are divided by $s_y^2$ so
    that residual terms are dimensionless and numerically comparable. All model
    diagnostics were computed on the original response scale; for the Neural
    ODE and Neural SDE this required transforming fitted normalized values back
    to the scale of $y_i^{[j]}$.